欢迎进入视觉组通识培训。

本卷不涉及装甲板识别、自瞄链路或哨兵导航，目标是使学员能够在 Linux 环境下使用现代 C++ 完成工程构建，并使用 git 进行版本管理。上述能力是后续 ROS2、OpenCV、深度学习等专项学习的共同基础。

教材分为三卷，建议按下列顺序学习：
\begin{enumerate}
    \item \textbf{卷一 · 基础通识}：Linux 环境、现代 C++、CMake
    \item \textbf{卷二 · 专项技术栈}：ROS2、图像处理、深度学习、估计与控制、导航等
    \item \textbf{卷三 · RM 视觉概述}：比赛场景下的系统结构与工程准则
\end{enumerate}

\textbf{本卷考核要求：}
\begin{enumerate}
    \item 独立完成可用的 Ubuntu 开发环境，或使用队伍统一镜像
    \item 掌握常用命令行与 git 基本工作流
    \item 能够编写类，理解 lambda、引用与移动语义，并使用智能指针与常用 STL
    \item 能够使用 CMake 编译多文件工程
\end{enumerate}

未通过本卷考核前，不宜直接修改实车业务仓库代码。配置环境前须完成系统备份。标 $*$ 的小节为选修，不计入通识考核。

工欲善其事，必先利其器。应先掌握工具与语言，再追求算法上限。
